\section{Carbon trong tự nhiên và nhiên liệu hóa thạch}

\subsection{Carbon tồn tại ở nhiều dạng}

Carbon có mặt trong khí quyển, sinh vật, đất đá và nhiên liệu. Một số dạng thường gặp:
\begin{itemize}
\item đơn chất carbon: than, graphite, diamond;
\item carbon dioxide: \ce{CO2};
\item muối carbonate: \ce{CaCO3}, \ce{MgCO3};
\item hợp chất hữu cơ trong cơ thể sinh vật và nhiên liệu.
\end{itemize}

\begin{ghinho}{Chu trình carbon}
Carbon liên tục chuyển từ môi trường sang sinh vật rồi trở lại môi trường. Có thể hình dung theo chuỗi:
\[
\boxed{\ce{CO2}\ \text{trong khí quyển}}
\]
\[
\Downarrow\ \text{quang hợp}
\]
\[
\boxed{\text{chất hữu cơ trong sinh vật}}
\]
\[
\Downarrow\ \text{hô hấp, phân hủy, đốt cháy}
\]
\[
\boxed{\ce{CO2}\ \text{trở lại môi trường}}.
\]
\end{ghinho}

Quang hợp làm giảm \ce{CO2} trong không khí, còn hô hấp, phân hủy và đốt cháy làm carbon trở lại khí quyển chủ yếu dưới dạng \ce{CO2}.

\subsection{Nhiên liệu hóa thạch}

Than, dầu mỏ và khí thiên nhiên được hình thành trong thời gian rất dài từ vật chất hữu cơ. Chúng là nguồn năng lượng quan trọng nhưng không phải nguồn tài nguyên có thể phục hồi nhanh.

Khi đốt nhiên liệu chứa carbon:
\[
\ce{C + O2 -> CO2}.
\]
Với hydrocarbon, carbon tạo \ce{CO2} và hydrogen tạo \ce{H2O} khi cháy hoàn toàn.

\begin{notechem}
Đốt nhiên liệu hóa thạch chuyển một lượng carbon đã được lưu giữ lâu dài trong lòng đất thành \ce{CO2} trong khí quyển. Vì vậy, sử dụng nhiên liệu hóa thạch với quy mô lớn làm tăng phát thải \ce{CO2}.
\end{notechem}

\subsection{Methane và khí nhà kính}

Methane \ce{CH4} là thành phần chính của khí thiên nhiên và cũng có thể sinh ra từ quá trình phân hủy chất hữu cơ trong điều kiện thiếu oxygen, chẳng hạn ở bãi chôn lấp, ruộng ngập nước hoặc hầm biogas.

Khi cháy hoàn toàn:
\[
\ce{CH4 + 2O2 -> CO2 + 2H2O}.
\]

Cả \ce{CO2} và \ce{CH4} đều góp phần vào hiệu ứng nhà kính. Do đó, giảm thất thoát methane và sử dụng nhiên liệu hợp lí đều có ý nghĩa đối với môi trường.

\begin{vidu}
Một trang trại thu khí biogas từ chất thải chăn nuôi để đun nấu thay cho việc để khí sinh ra thoát trực tiếp vào không khí. Nêu hai lợi ích chính của việc làm này.
\end{vidu}
\begin{loigiai}
Biogas được tận dụng làm nhiên liệu nên giảm nhu cầu sử dụng một phần nhiên liệu khác. Đồng thời methane sinh ra được thu gom và đốt cháy có kiểm soát thay vì thoát trực tiếp ra môi trường.

\textbf{Vậy,} việc thu biogas vừa tận dụng được nguồn nhiên liệu, vừa hạn chế phát tán methane trực tiếp vào khí quyển.
\end{loigiai}

\begin{tuluyen}
Phân loại các quá trình sau thành quá trình làm giảm hoặc làm tăng lượng \ce{CO2} trong khí quyển:
\begin{tasks}(2)
\task quang hợp của cây xanh;
\task đốt than;
\task hô hấp của sinh vật;
\task đốt khí thiên nhiên.
\end{tasks}
\dapso{Giảm: quang hợp. Tăng: đốt than, hô hấp, đốt khí thiên nhiên.}
\end{tuluyen}
